\begin{table}[t]
\centering
\caption{Why the null result is informative in the confirmatory C100 run ($n_{\text{pairs}}=1193$).}
\label{tab:negative-result-diagnostics}
\begin{tabular}{p{0.33\linewidth}p{0.61\linewidth}}
\toprule
Diagnostic dimension & Interpretation \\
\midrule
Effect-size detectability & $\Delta$ nDCG@10 ($-9.2588\times10^{-6}$) is near zero and far below the practical threshold (0.02), indicating no detectable material gain. \\
Confidence interval meaning & Uncertainty is compatible with negligible effects around zero; intervals do not support a superiority narrative under the confirmatory protocol. \\
Practical-threshold comparison & The observed magnitude fails the operational relevance criterion by orders of magnitude. \\
Latency penalty interpretation & Kalman/Mean latency ratio $=1.0723$ implies about $7.23\%$ runtime overhead at matched FLOPs (ratio $=1.0$). \\
Baseline competitiveness & Mean fusion remains a strong baseline because no primary or secondary quality uplift is demonstrated. \\
Uncertainty calibration impact & The tested uncertainty weighting does not convert calibration information into measurable retrieval gains in this slice. \\
Evidence needed to reverse conclusion & Future confirmatory evidence must show positive direction, practical effect, Holm-adjusted significance, stable non-negligible intervals, and independent replication. \\
\bottomrule
\end{tabular}
\end{table}
